\documentclass[UTF8,openany,oneside,zihao=-4]{ctexbook}

\usepackage[a4paper,margin=2cm,headheight=14pt]{geometry}
\usepackage{amsmath,amssymb}
\usepackage{xcolor}
\usepackage{enumitem}
\usepackage{fancyhdr}
\usepackage{listings}
\usepackage{titlesec}
\usepackage{tcolorbox}
\tcbuselibrary{skins,breakable}

% ========== 颜色 ==========
\definecolor{codeblue}{HTML}{0550AE}
\definecolor{codegray}{HTML}{6E7781}
\definecolor{codebg}{HTML}{F6F8FA}
\definecolor{warnred}{HTML}{CF222E}
\definecolor{fixgreen}{HTML}{116329}
\definecolor{teal}{HTML}{0E6E6E}

% ========== 页眉页脚 ==========
\pagestyle{fancy}
\fancyhf{}
\fancyhead[L]{\small\textbf{错题本} | WA 排查}
\fancyhead[R]{\small\thepage}
\renewcommand{\headrulewidth}{0.3pt}

% ========== 段落与列表 ==========
\setlength{\parindent}{0pt}
\setlength{\parskip}{0.4em}
\setlist[itemize]{leftmargin=1.8em,itemsep=0.15em,topsep=0.15em}
\setlist[enumerate]{leftmargin=2em,itemsep=0.2em,topsep=0.2em}

% ========== 标题样式 ==========
\titleformat{\chapter}{\LARGE\bfseries}{\thechapter}{0.7em}{}[\vspace{0.2em}\hrule]
\titleformat{\section}{\large\bfseries\color{codeblue}}{\thesection}{0.5em}{}
\titleformat{\subsection}{\normalsize\bfseries\color{teal}}{\thesubsection}{0.5em}{}

% ========== 代码样式 ==========
\lstdefinestyle{cpp}{
  language=C++,
  basicstyle=\ttfamily\small,
  keywordstyle=\color{codeblue}\bfseries,
  commentstyle=\color{codegray},
  stringstyle=\color{warnred},
  backgroundcolor=\color{codebg},
  frame=single,
  rulecolor=\color{codebg},
  breaklines=true,
  columns=fullflexible,
  keepspaces=true,
  showstringspaces=false,
  tabsize=4
}
\lstset{style=cpp}

% ========== 提示框 ==========
\newtcolorbox{warn}{
  colback=red!3,colframe=warnred,
  fonttitle=\bfseries,title={\color{warnred}坑},
  boxrule=0.4pt,left=4pt,right=4pt,top=2pt,bottom=2pt,
  breakable
}
\newtcolorbox{tip}{
  colback=green!3,colframe=fixgreen,
  fonttitle=\bfseries,title={\color{fixgreen}解法},
  boxrule=0.4pt,left=4pt,right=4pt,top=2pt,bottom=2pt,
  breakable
}
\newtcolorbox{note}{
  colback=blue!3,colframe=codeblue,
  boxrule=0.4pt,left=4pt,right=4pt,top=2pt,bottom=2pt,
  breakable
}

\begin{document}

% ========== 封面 ==========
\begin{titlepage}
\centering
\vspace*{3cm}
{\Huge\bfseries 错题本\par}
\vspace{0.6cm}
{\Large WA / RE / TLE 排查速查手册\par}
\vspace{1.5cm}
{\normalsize 打印用 · 比赛翻阅 · 提交前过一遍\par}
\vfill
{\large \today\par}
\end{titlepage}

\tableofcontents
\clearpage

% ==========================================================
\chapter{板子}
% ==========================================================

比赛开局直接粘贴的模板。注意 \texttt{\#define int long long} 之后 main 必须写 \texttt{signed main}。

\begin{lstlisting}
#include <bits/stdc++.h>
using namespace std;

#define int long long
#define endl '\n'

const int inf = 2e18;
const int N = 1e6 + 10;

void solve() {
    int n;
    cin >> n;

}

signed main() {
    ios::sync_with_stdio(0);
    cin.tie(0);
    int t = 1;
    cin >> t;
    while (t--) {
        solve();
    }
    return 0;
}
\end{lstlisting}

% ==========================================================
\chapter{提交前速查}
% ==========================================================

\begin{note}
大多数 WA 不是算法错了，是细节把正确思路拖下水。提交前花 30 秒过一遍这页。
\end{note}

\begin{enumerate}
  \item 有没有 \lstinline|1 << i| 没写成 \lstinline|1LL << i|？
  \item \lstinline|inf + x| 会不会爆？转移前加 \lstinline|if(dist[u] != inf)|。
  \item 两个大数相乘有没有步步取模？
  \item 多测的数组/图清空了吗？\lstinline|solve()| 里有没有提前 \lstinline|return| 导致输入没读完？
  \item 减法取模补 \lstinline|+ MOD| 了吗？除法用逆元了吗？
  \item $n=0,1,2$、全相同、全负数、无解、不连通测过了吗？
  \item 输出大小写对吗？下标是 0-based 还是 1-based？行末空格？
\end{enumerate}

% ==========================================================
\chapter{通用 WA 错误}
% ==========================================================

\section{溢出}

\begin{warn}
\begin{itemize}
  \item \texttt{1 << i}：常数 1 是 32 位，$i \ge 31$ 就炸。写 \texttt{1LL << i}。
  \item \texttt{inf + x}：$2 \times 10^{18} + x$ 爆 long long 变负数。松弛/转移前判不可达。
  \item 大数乘法：$10^{18} \times 10^{18}$ 直接爆。用 \texttt{\_\_int128} 或步步取模。
  \item \texttt{pow(5,2)} 可能返回 24.999...，强转 int 变 24。整数幂用快速幂。
\end{itemize}
\end{warn}

\section{多测}

\begin{warn}
\begin{itemize}
  \item \textbf{串流}：\lstinline|solve()| 里提前 return，这组数据没读完，后面全乱。
  \item \textbf{残留}：上一组 $n=10$，这组 $n=5$，只清 $1 \sim 5$ 留下 $6 \sim 10$ 的脏数据。
  \item \textbf{重名}：全局 \lstinline|int n;| 和局部 \lstinline|int n;| 同时存在，DFS 读到的是全局的 0。
\end{itemize}
\end{warn}

\section{取模}

\begin{itemize}
  \item 减法：\lstinline|(a - b % MOD + MOD) % MOD|，防负数。
  \item 除法：不能直接除，要乘逆元 \lstinline|a * qpow(b, MOD-2) % MOD|。
  \item 连乘：\lstinline|A * B % MOD * C % MOD|，步步取模防中间爆。
\end{itemize}

\section{边界}

\begin{itemize}
  \item $n=0,1,2$ 特判。
  \item 全相同、全为 $10^9$、全为负数。
  \item 图：不连通、链、菊花、自环、重边。
  \item 无解时输出什么？$-1$？$0$？\texttt{Impossible}？
\end{itemize}

\section{输出格式}

\begin{itemize}
  \item YES/Yes/yes 大小写。
  \item 下标 +1 还是不 +1。
  \item 行末空格、最后换行。
  \item 多解时要不要字典序最小。
\end{itemize}

\section{杂项}

\begin{itemize}
  \item \textbf{贪心假了}：无法证明就怀疑，考虑 DP / 二分答案。
  \item \textbf{比较器}：\lstinline|sort| 的 cmp 用了 \lstinline|<=| 会破坏严格弱序，必须用 \lstinline|<|。
  \item \textbf{未初始化}：本地对但提交 WA，检查局部变量有没有赋初值。
  \item \textbf{哈希被卡}：\lstinline|unordered_map| 在 CF 上容易被 hack，加 custom hash 或换 \lstinline|map|。
\end{itemize}

\section{15 分钟还没找到 WA}

\begin{enumerate}
  \item 停止看代码。手造小样例（含负数、0、1、乱序）。
  \item 纸上跑一遍算法。
  \item 写暴力（$O(n^3)$ 也行），写随机数据生成器，对拍。
  \item 找到最小错误数据，定位第一处不一致。
\end{enumerate}

% ==========================================================
\chapter{错题复盘与最小反例}
% ==========================================================

这一章不是再罗列“某某算法会错”，而是把一次错误真正变成下一次不会再犯的检查项。每道题至少留下四样东西：

\begin{enumerate}
  \item 最小错误输入：尽量短，但必须仍然能复现错误；
  \item 错误原因：指出哪一个不变量、边界或题意被破坏；
  \item 修复后的关键代码：只保留真正改变结果的几行；
  \item 回归测试：以后每次改模板都重新运行，防止旧错误回来。
\end{enumerate}

\section{一次 WA 的完整复盘流程}

\begin{enumerate}
  \item \textbf{先确认输入输出没有错：} 样例是否完整复制，是否漏读了 $T$、多组边、空行或字符串；输出大小写、空格和换行是否符合题面。
  \item \textbf{再缩小数据：} 删除无关点、边、操作，优先保留第一个产生错误的操作。一个能在 4 个数上复现的错误，不要拿 4 万个数解释。
  \item \textbf{写出不变量：} 例如堆顶是当前最短距离、线段树节点的和对应整个区间、动态规划状态已经处理完前缀、回滚前后的历史栈长度相同。
  \item \textbf{对照错误发生前后：} 在第一次不符合不变量的位置打印变量，而不是只打印最终答案。
  \item \textbf{写进回归样例：} 把这个最小输入和正确输出放入模板旁边，下一次修改代码时必须重新跑。
\end{enumerate}

\begin{warn}
不要一看到 WA 就重写整份代码。先找到“第一个错误状态”，否则新代码即使过了当前样例，也不知道究竟修复了什么。
\end{warn}

\section{多测与 \texttt{void solve()}：最容易被忽略的真实错误}

\subsection{中途 return 会把下一组输入读乱}

下面的代码在无解时直接返回，当前测试剩余的边仍在输入流中，下一组的 $n,m$ 就会读成上一组的边：

\begin{lstlisting}
void solve() {
    int n, m;
    cin >> n >> m;
    for (int i = 0; i < m; ++i) {
        int u, v;
        cin >> u >> v;
        if (u == v) return; // 错：剩余输入没有读完
    }
    cout << "YES\n";
}
\end{lstlisting}

如果输入是：

\begin{lstlisting}
2
3 2
1 1
1 2
1 0
\end{lstlisting}

第一组提前返回后，第二组的 \lstinline|n,m| 可能被读成 \lstinline|1,2|，错误会看起来像“第二组算法随机坏了”。正确写法是读完当前组，再用标志记录非法状态：

\begin{lstlisting}
void solve() {
    int n, m;
    cin >> n >> m;
    bool ok = true;
    for (int i = 0; i < m; ++i) {
        int u, v;
        cin >> u >> v;
        if (u == v) ok = false;
    }
    cout << (ok ? "YES" : "NO") << '\\n';
}
\end{lstlisting}

\subsection{全局容器和节点池必须按题目边界清理}

多测最小测试：第一组规模大，第二组规模小。它可以专门暴露“只清理前 $n$ 个位置”或“节点编号没有归零”：

\begin{lstlisting}
2
5
1 2 3 4 5
1
9
\end{lstlisting}

如果第二组的答案仍然包含第一组的数据，优先检查：

\begin{itemize}
  \item \lstinline|vector| 是否调用了 \lstinline|clear()| 或重新赋值；
  \item 图的边表、\lstinline|head|、\lstinline|edge_count| 是否一起清空；
  \item 字典树、主席树、回文自动机的 \lstinline|tot| 是否恢复到根节点之后；
  \item 线段树的 \lstinline|sum|、\lstinline|lazy|、\lstinline|tag| 是否全部重新 \lstinline|build|；
  \item 时间戳数组是否使用了新的 \lstinline|stamp|，而不是只清空一半下标。
\end{itemize}

\subsection{多测回归样例的固定格式}

以后给任何板子补样例时，优先使用下面的四组结构：

\begin{enumerate}
  \item $n=1$ 或空结构：检查根、叶子、空集合和初始化；
  \item 退化结构：链、星、全相同、全负、无边或单条边；
  \item 正常结构：至少有两条竞争路径、两种状态或两个不同区间；
  \item 两组连续测试：第一组故意较大，第二组很小，检查状态污染。
\end{enumerate}

\section{边界样例设计 SOP}

\subsection{先列“变化维度”，再写数据}

不要随机生成一堆看起来很大的数。先列出模板真正依赖的变化维度：

\begin{itemize}
  \item 数组：长度 $0,1,2$、重复值、负数、最大值、单调和逆序；
  \item 区间：单点、整段、相邻、不相交、完全包含、端点重合；
  \item 图：无边、链、星、环、重边、自环、不连通、负边；
  \item 树：单点、链、星、平衡树、根是答案、根不在查询路径；
  \item 字符串：空串、单字符、全相同、周期串、重叠匹配、无匹配；
  \item 几何：重合、相切、两点交、平行、共线、退化多边形；
  \item 多测：大组后小组、空组后正常组、节点池复用。
\end{itemize}

\subsection{每个大型板子至少四组核心测试}

建议每个超过十行的模板都保留以下四组，而不是只保留随机样例：

\begin{center}
\begin{tabular}{p{0.20\textwidth}p{0.33\textwidth}p{0.32\textwidth}}
\hline
类型 & 目标 & 典型输入 \tabularnewline
\hline
边界一 & 空结构或最小规模 & $n=0,1$、单点、空边集 \tabularnewline
边界二 & 退化或等号分支 & 全相同、相切、重复边、$x=\text{second\_max}$ \tabularnewline
正常一 & 核心转移真的发生 & 两条竞争路径、重叠区间、至少一次旋转/合并 \tabularnewline
正常二 & 多阶段组合 & 先修改再查询、先连边再断边、多个版本交叉查询 \tabularnewline
\hline
\end{tabular}
\end{center}

如果板子存在随机化或浮点误差，再增加一组“同答案不同顺序”的测试：输入随机打乱后，答案允许顺序不同，但数值和不变量必须相同。

\section{最小反例库：按错误类型直接查}

\subsection{下标与半开区间}

所有只处理 \lstinline|[l,r)| 的代码都要测试 \lstinline|l=0,r=1|、\lstinline|l=0,r=n| 和相邻区间：

\begin{lstlisting}
数组：      [10, 20, 30]
区间一：    [0, 1)  只包含 10
区间二：    [1, 3)  包含 20,30
区间三：    [0, 3)  包含全部元素
区间四：    [1, 2)  只包含 20
\end{lstlisting}

如果区间三答案少一个元素，通常是把右端点当成闭区间；如果区间一访问到 \lstinline|a[1]|，通常是把右端点错误地包含进来了。

\subsection{前缀和与差分}

对数组 \lstinline|[1,2,3]| 依次执行 \lstinline|add(1,1,5)|、\lstinline|add(2,3,-2)|，最终应为 \lstinline|[6,0,1]|。这个样例同时检查：

\begin{itemize}
  \item 单点更新是否正确；
  \item 负数更新是否正确；
  \item \lstinline|r+1| 是否写在正确位置；
  \item 查询 \lstinline|[1,3]| 是否为 $7$，而不是只返回最后一次更新。
\end{itemize}

\subsection{图论：不可达、重边、自环}

\begin{lstlisting}
4
1 0
2 1 1 2
2 2 1 2 1 2
3 3 1 1 1 2 2 3
\end{lstlisting}

四组分别测试单点、单边、重边和自环。不要默认题目没有自环；如果算法是桥、割点、强连通分量或拓扑排序，自环的处理完全不同：自环会让有向图出现环，但不会把无向单点自动变成桥。

\subsection{最短路：旧堆项和不可达点}

\begin{lstlisting}
3 3
1 2 10
1 3 1
3 2 1
\end{lstlisting}

从 1 到 2 的答案是 2。Dijkstra 会先把 2 以 10 放进堆，再通过 3 改成 2；弹出旧的 \lstinline|(10,2)| 时必须跳过。再增加一个没有任何入边的点，检查不可达点是否保持 \lstinline|INF|，而不是参与 \lstinline|INF+w| 的溢出运算。

\subsection{动态规划：全负、空转移和答案位置}

最大子段和的回归样例：

\begin{lstlisting}
4
1
-5
3
1 -2 1
4
-3 -2 -5 -1
5
2 -1 2 -1 2
\end{lstlisting}

答案分别是 \lstinline|-5,1,-1,4|。如果把 \lstinline|dp| 初值写成 0，第一组就会错误；如果只在转移时更新答案而不检查最后一个位置，第四组也容易错。

\subsection{线段树：整段标记后查询子区间}

初始数组 \lstinline|[1,1,1,1]|，先对 \lstinline|[1,4]| 加 5，再查询 \lstinline|[2,3]|，答案必须为 12；随后对 \lstinline|[2,2]| 加 -3，再查整段，答案为 18。这个顺序专门验证 \lstinline|push_down|：只查整段时懒标记可能一直藏着，只有访问子区间才会暴露错误。

\subsection{字符串：重叠和周期}

文本 \lstinline|abababa|、模式 \lstinline|aba| 的匹配起点为 \lstinline|0,2,4|，总数为 3；文本 \lstinline|aaaaa|、模式 \lstinline|aaa| 的总数也为 3。匹配成功后不能把模式指针直接归零，必须回到前缀函数给出的最长可复用前缀。

\subsection{计算几何：相切、共线和重复点}

\begin{itemize}
  \item 两个半径为 1 的圆，圆心距 2：一个交点；圆心距 1：两个交点；圆心距 3：无交点。
  \item 三点 \lstinline|(0,0),(1,0),(3,0)|：外接圆应退化成最远两点的直径圆，不能除以零。
  \item 最近点对中加入两个相同点：答案必须为 0，不能用一个正的初始答案掩盖重复点。
\end{itemize}

\section{错误代码与修复代码对照}

\subsection{负数取模}

\begin{lstlisting}
// 错误：C++ 的负数取模结果仍可能为负
long long add = (a - b) % MOD;

// 修复：先把减数归一化，再整体加模
long long add = (a - b % MOD + MOD) % MOD;
\end{lstlisting}

回归输入：\lstinline|a=3,b=5,MOD=7|，答案应为 5，而不是 C++ 直接取模得到的 -2。

\subsection{区间最小值的错误懒标记}

\begin{lstlisting}
// 错误：把 chmin 当成加法，直接给整个区间加 tag
tree[p].mx = min(tree[p].mx, x);
tree[p].sum += x * len; // 完全错误

// 修复思路：只有当前最大值会下降，
// 必须维护最大值、次大值、最大值出现次数和区间和。
if (second_mx[p] < x && x < mx[p]) {
    sum[p] -= (mx[p] - x) * cnt_mx[p];
    mx[p] = x;
}
\end{lstlisting}

回归数组 \lstinline|[1,3,3,7]|：执行整段 \lstinline|chmin(5)| 后和为 12，再执行 \lstinline|[2,3] chmin(2)| 后该区间和为 4。第二次操作不能只看根节点最大值，必须按次大值条件决定是否继续递归。

\subsection{LCA 根节点初始化}

\begin{lstlisting}
// 错误：根的父亲没有明确设为 0，倍增表含有旧数据
dfs(root, parent[root]);

// 修复：每组数据先固定根和所有根祖先
parent[root] = 0;
depth[root] = 0;
for (int j = 0; j < LOG; ++j) up[root][j] = 0;
dfs(root, 0);
\end{lstlisting}

回归树：单点、链 \lstinline|1-2-3-4|、星 \lstinline|1-{2,3,4}|，分别查询 \lstinline|lca(1,1)|、\lstinline|lca(2,4)|、\lstinline|lca(3,4)|。根表没有清空时，单点组往往会错误继承上一组根的祖先。

\section{可直接改名使用的对拍框架}

下面的框架不替代暴力解，它只负责反复生成小数据、运行标准程序和暴力程序，并在第一次不一致时停下。把三个可执行文件名替换成自己的文件即可：

\begin{lstlisting}
#include <bits/stdc++.h>
using namespace std;

int main() {
    mt19937 rng(chrono::steady_clock::now().time_since_epoch().count());
    for (int tc = 1; tc <= 100000; ++tc) {
        // generator.exe 负责把一组随机数据写入 input.txt
        system("generator.exe > input.txt");
        system("std.exe < input.txt > std.out");
        system("brute.exe < input.txt > brute.out");
        ifstream a("std.out"), b("brute.out");
        string x, y;
        bool same = true;
        while (a >> x || b >> y) {
            if (x != y) { same = false; break; }
            x.clear(); y.clear();
        }
        if (!same) {
            cerr << "Mismatch on test " << tc << '\\n';
            cerr << "Input is saved in input.txt\\n";
            return 0;
        }
    }
    cerr << "All tests passed\\n";
}
\end{lstlisting}

对拍生成器不要一开始就追求大数据。先覆盖：$n=1$、全相同、完全逆序、只有一条边、重复边、无解、答案在端点、随机小数据。暴力程序必须比标准程序更简单，即使慢到 $O(n^3)$ 也可以；对拍的目标是正确，不是速度。

\section{错题卡：以后每道题只记录一页}

\begin{center}
\begin{tabular}{p{0.25\textwidth}p{0.62\textwidth}}
\hline
字段 & 记录内容 \tabularnewline
\hline
题目与提交时间 & 题号、语言、提交结果、使用的模板 \tabularnewline
最小错误输入 & 删除哪些元素后仍然能复现 \tabularnewline
错误类别 & 题意、边界、下标、溢出、状态、清空、复杂度、输出 \tabularnewline
被破坏的不变量 & 例如“当前区间和应等于四个子区间之和” \tabularnewline
错误代码 & 只贴最小的 3--10 行，不要整页复制 \tabularnewline
修复动作 & 改了哪一个条件、初始化或数据结构 \tabularnewline
回归样例 & 输入、正确输出、以后什么时候重跑 \tabularnewline
\hline
\end{tabular}
\end{center}

\begin{tip}
真正有价值的错题记录不是“这题用了线段树”，而是“我把整段懒标记下传前直接访问了子区间；最小反例是先改全区间再查中间一格”。后者可以迁移到下一道题，前者只能记住一道题。
\end{tip}

\section{提交前五分钟固定动作}

\begin{enumerate}
  \item 用题面约束重新确认 \lstinline|int|、\lstinline|long long|、\lstinline|__int128| 和浮点类型；
  \item 手算一组 $n=1$，一组重复值/重边，一组无解或不可达；
  \item 检查 \lstinline|solve()| 是否完整读完当前测试，所有全局状态是否清空；
  \item 检查所有区间端点是闭区间还是半开区间，下标是 0-based 还是 1-based；
  \item 检查输出顺序、大小写、精度和无解格式；
  \item 如果仍然不放心，立刻写暴力和对拍，不要继续凭感觉改条件。
\end{enumerate}

% ==========================================================
\chapter{通用 RE 错误}
% ==========================================================

\begin{note}
RE 不像 WA 有输出可以对比——程序直接崩了，没有任何线索。
排查思路：先定位\textbf{崩在哪一行}，再判断\textbf{哪类原因}。
\end{note}

\section{RE 排查流程}

\begin{enumerate}
  \item \textbf{本地复现}：把样例贴进来跑，看能不能在本地炸。
  \item \textbf{加输出定位}：在关键位置加 \lstinline|cerr << "checkpoint 1" << endl;|，逐步缩小崩溃区间。
  \item \textbf{极端数据}：$n=0$、$n=1$、$n=N_{\max}$、全相同、全极值，单独测。
  \item \textbf{开编译选项}：本地编译加 \lstinline|-fsanitize=address,undefined -g|，能直接报出越界和 UB 的行号。
  \item \textbf{对照常见原因清单}：逐条过下面的列表。
\end{enumerate}

\section{数组越界}

\begin{warn}
\begin{itemize}
  \item \textbf{开小了}：$N = 10^5 + 10$ 但输入 $n$ 可以到 $2 \times 10^5$。仔细看数据范围再开数组。
  \item \textbf{下标为负}：\lstinline|a[i - k]| 当 $i < k$ 时变负数。用前先判 \lstinline|if(i >= k)|。
  \item \textbf{下标为 0}：1-based 存图但 \lstinline|head[0]| 没初始化；或反过来 0-based 但循环从 1 开始漏了 \lstinline|a[0]|。
  \item \textbf{邻接表边数}：无向图每条边存两次，边数组至少开 $2M$。
  \item \textbf{线段树}：节点数组至少 $4N$（不是 $2N$）。
  \item \textbf{组合数/阶乘表}：预处理到 $N$ 但查询 $C(n, m)$ 时 $n$ 可能到 $2N$。
\end{itemize}
\end{warn}

\section{栈溢出（递归过深）}

\begin{warn}
\begin{itemize}
  \item \textbf{链状图 DFS}：$n = 10^5$ 的链递归深度 $10^5$ 层，默认栈（1\,MB / 8\,MB）直接爆。
  \item \textbf{解决方案}：
    \begin{itemize}
      \item 改用 BFS / 手写栈模拟 DFS。
      \item Linux 下加 \lstinline|ulimit -s unlimited|（某些 OJ 支持）。
      \item Windows 下编译加 \lstinline|-Wl,--stack,268435456|（256\,MB 栈）。
    \end{itemize}
  \item \textbf{无限递归}：忘记写 \lstinline|return| / base case，或者 DFS 忘标 \lstinline|vis[u] = true|，同一个节点反复进入。
  \item \textbf{树形 DP 忘跳父亲}：\lstinline|dfs(v, u)| 里没写 \lstinline|if(v == fa) continue|，父子互相递归。
\end{itemize}
\end{warn}

\section{除以零 / 取模零}

\begin{itemize}
  \item \lstinline|a / b| 或 \lstinline|a % b| 当 $b = 0$ 时直接 RE。
  \item 常见场景：二分 \lstinline|mid = (l + r) / 2| 本身不会为零，但 \lstinline|n / mid| 会在 \lstinline|mid == 0| 时炸。
  \item GCD 循环里 \lstinline|a % b|，如果 $b$ 初始就是 0 则第一轮就崩。
  \item \textbf{解法}：所有除法/取模前加 \lstinline|if(b == 0)| 特判，或保证逻辑上 $b \ne 0$。
\end{itemize}

\section{非法内存访问}

\begin{itemize}
  \item \textbf{空指针}：链表/树结构访问 \lstinline|nullptr->val|。
  \item \textbf{野指针}：\lstinline|delete| 后继续使用。竞赛中少用动态内存，优先用静态数组 + 内存池。
  \item \textbf{VLA 过大}：\lstinline|int a[n]| 在栈上分配，$n = 10^6$ 就可能爆栈。改用全局数组或 \lstinline|vector|。
\end{itemize}

\section{STL 误用}

\begin{warn}
\begin{itemize}
  \item \textbf{空容器取元素}：
    \begin{itemize}
      \item \lstinline|v.back()| / \lstinline|v.front()| 在 \lstinline|v.empty()| 时 UB/RE。
      \item \lstinline|s.top()| / \lstinline|q.front()| 在空栈/空队列时崩。
      \item 用前必须判 \lstinline|!v.empty()|。
    \end{itemize}
  \item \textbf{迭代器失效}：
    \begin{itemize}
      \item \lstinline|vector| 的 \lstinline|push_back| 可能导致扩容，所有迭代器和指针失效。
      \item \lstinline|set/map| 的 \lstinline|erase| 后原迭代器失效，用 \lstinline|it = s.erase(it)| 或先 \lstinline|next(it)|。
    \end{itemize}
  \item \textbf{越界访问}：\lstinline|v[i]| 不做边界检查（\lstinline|v.at(i)| 会抛异常）。
  \item \textbf{比较器不满足严格弱序}：\lstinline|sort| 的 cmp 用 \lstinline|<=| 会导致内部断言失败 → RE（不只是 WA！）。
\end{itemize}
\end{warn}

\section{整数未定义行为}

\begin{itemize}
  \item \textbf{有符号溢出}：C++ 中有符号整数溢出是 UB，编译器可能生成崩溃代码。
  \item \textbf{移位越界}：\lstinline|1 << 31| 对 32 位 int 是 UB（符号位）；\lstinline|1 << 64| 对 64 位也是 UB。
  \item \textbf{负数移位}：\lstinline|(-1) << k| 和 \lstinline|x >> 负数| 都是 UB。
  \item 开 \lstinline|-fsanitize=undefined| 可以在本地立刻暴露这些问题。
\end{itemize}

\section{常见 RE 清单速查}

\begin{note}
提交 RE 后，按此表逐条排查：
\end{note}

\begin{enumerate}
  \item 数组大小够吗？下标有没有负数或超上界？
  \item 递归深度超过 $10^4$ 了吗？有没有无限递归？
  \item 有除法或取模吗？除数可能为 0 吗？
  \item 有 \lstinline|sort| 自定义比较器吗？是否满足严格弱序（\lstinline|<| 而非 \lstinline|<=|）？
  \item 用了 \lstinline|stack/queue/vector| 的 \lstinline|top()/front()/back()| 吗？容器可能为空吗？
  \item 用了 \lstinline|1 << k| 吗？$k \ge 31$ 时要写 \lstinline|1LL << k|。
  \item 有 VLA（\lstinline|int a[n]|）吗？$n$ 大时改用全局数组。
  \item 线段树开了 $4N$ 吗？并查集/倍增表第二维开够了吗？
  \item 多测时上一组残留的大数组会不会越界访问？
  \item 本地加 \lstinline|-fsanitize=address,undefined -g| 编译，跑一遍样例。
\end{enumerate}

\begin{tip}
RE 和 WA 的本质区别：WA 是逻辑错误，RE 是\textbf{访问了不该访问的东西}。
绝大多数 RE 都是数组越界和栈溢出，优先查这两个。
\end{tip}

% ==========================================================
\chapter{图论与树}
% ==========================================================

\section{建图}

\begin{itemize}
  \item 无向图边数开 $2M$（正反各一条）。
  \item 点属性开 $N$，边属性开 $M$，别搞混。
  \item 多测清空要按上一组的范围清，链式前向星 \lstinline|cnt = 0|。
  \item 编号 0-based 还是 1-based，循环和起点要一致。
\end{itemize}

\section{特殊图}

\begin{itemize}
  \item \textbf{不连通}：没保证连通就外层 \lstinline|for(i=1;i<=n;i++) if(!vis[i]) dfs(i)|。
  \item \textbf{重边}：邻接矩阵 \lstinline|g[u][v] = min(g[u][v], w)|。
  \item \textbf{自环}：影响入度、度数、连通性。必要时 \lstinline|if(u==v) continue|。
  \item \textbf{链}：递归深搜会爆栈。\textbf{菊花}：暴力会退化成 $O(n^2)$。
\end{itemize}

\section{DFS / BFS}

\begin{itemize}
  \item 无向树 DFS 传父亲：\lstinline|dfs(u, fa)|，跳过 \lstinline|if(v==fa) continue|。
  \item BFS 入队时立刻标 \lstinline|vis[v]=true|，不要等出队。
  \item 找所有路径要回溯：离开前 \lstinline|vis[u]=0|。
\end{itemize}

\section{最短路}

\begin{itemize}
  \item Dijkstra 堆优化：出队判 \lstinline|if(vis[u]) continue|。
  \item 有负权 $\Rightarrow$ Dijkstra 失效，用 Bellman-Ford。
  \item Floyd：初始化 \lstinline|inf|，对角线 0，循环顺序 $k,i,j$。
  \item 松弛防溢出：\lstinline|if(dist[u]!=inf && dist[v]>dist[u]+w)|。
\end{itemize}

\section{拓扑排序}

\begin{itemize}
  \item 跑完检查 \lstinline|cnt==n|，不等说明有环。
  \item 要字典序最小 $\Rightarrow$ 用小根堆代替普通队列。
\end{itemize}

\section{LCA / 树形 DP}

\begin{itemize}
  \item 倍增第二维开 $\lceil\log_2 N\rceil$，预处理外层枚举步长、内层枚举节点。
  \item 树形 DP：先 \lstinline|dfs(v)| 再用 \lstinline|dp[v]| 更新 \lstinline|dp[u]|（后序）。
  \item 多测时在 \lstinline|dfs(u)| 开头清当前点的 DP 状态。
\end{itemize}

\begin{tip}
图论一直 WA/RE 时：在测试数据里加孤立点、自环、重边、退化成链/菊花，很多隐藏 bug 会立刻暴露。
\end{tip}

\section{图论模型先判型}

拿到图论题先不要急着敲模板，先把题面翻译成下面四个问题：

\begin{enumerate}
  \item 边有没有方向？无向边必须加入两次；有向边不能反向乱加。
  \item 权值范围是什么？只有 $0/1$ 用 0-1 BFS，非负任意权值用 Dijkstra，有负权考虑 Bellman-Ford 或差分约束。
  \item 题目问的是一条路径、所有点的可达性、连通块，还是删点删边之后的结构？对应的算法完全不同。
  \item 图是否保证连通、无环、没有重边？题面没有保证的性质都不能默认。
\end{enumerate}

\begin{note}
 最容易发生的误判是“看到图就写 DFS”。DFS 只能完成遍历，不能自动解决最短路、拓扑序、强连通分量或割点问题。先写出状态含义：例如 \lstinline|dist[u]| 是源点到 $u$ 的最短距离，\lstinline|vis[u]| 只是“是否已经发现”，两者不是一回事。
\end{note}

\section{连通块、0-1 BFS 与最小生成树}

\subsection{连通块：外层循环不能漏}

只从 1 号点 DFS 只能得到 1 号点所在的连通块。要统计整个无向图的连通块，必须对所有点补一层外循环：

\begin{lstlisting}
vector<vector<int>> g;
vector<int> vis;

void dfs(int u) {
    vis[u] = 1;                    // 先标记，避免无向边来回递归
    for (int v : g[u]) {
        if (!vis[v]) dfs(v);
    }
}

int count_components(int n) {
    int components = 0;
    for (int u = 1; u <= n; ++u) {
        if (!vis[u]) {
            ++components;
            dfs(u);                // 每次启动 DFS 就找到一个新连通块
        }
    }
    return components;
}
\end{lstlisting}

\begin{warn}
 无向图的 DFS 如果只写 \lstinline|if (v == fa) continue|，遇到重边时仍可能出错：两条平行边中只有一条是“父边”，另一条不能被跳过。需要区分边编号时传 \lstinline|parent_edge|，不要只传父节点。
\end{warn}

\subsection{0-1 BFS：权值必须真的只有 0 和 1}

边权为 0 或 1 时，双端队列可以代替优先队列。松弛后，走 0 边放到队首，走 1 边放到队尾。它的关键不在“用了 deque”，而在于队列顺序保证较小距离优先被处理：

\begin{lstlisting}
struct Edge { int to, w; };       // w 必须是 0 或 1
const long long INF = (1LL << 60);

vector<long long> zero_one_bfs(int s, const vector<vector<Edge>>& g) {
    int n = (int)g.size() - 1;
    vector<long long> dist(n + 1, INF);
    deque<int> q;
    dist[s] = 0;
    q.push_front(s);

    while (!q.empty()) {
        int u = q.front();
        q.pop_front();
        for (auto e : g[u]) {
            if (dist[e.to] > dist[u] + e.w) {
                dist[e.to] = dist[u] + e.w;
                if (e.w == 0) q.push_front(e.to);
                else q.push_back(e.to);
            }
        }
    }
    return dist;
}
\end{lstlisting}

检查方式：把一条 0 边和一条 1 边交错放进图里，确认算法不会把“边数最少”误当成“权值和最小”；再故意放入权值 2，确认自己知道这时必须换 Dijkstra。

\subsection{Kruskal：最小生成树的三个返回值}

Kruskal 不只要返回总权值，还经常要返回选了多少条边。若最后选边数不是 $n-1$，图不连通，不能把当前总和当作一棵生成树：

\begin{lstlisting}
struct DSU {
    vector<int> fa, sz;
    DSU(int n) : fa(n + 1), sz(n + 1, 1) {
        iota(fa.begin(), fa.end(), 0);
    }
    int find(int x) { return fa[x] == x ? x : fa[x] = find(fa[x]); }
    bool unite(int x, int y) {
        x = find(x); y = find(y);
        if (x == y) return false;   // 会形成环，不能选
        if (sz[x] < sz[y]) swap(x, y);
        fa[y] = x; sz[x] += sz[y];
        return true;
    }
};

long long kruskal(int n, vector<tuple<int,int,int>> edges) {
    sort(edges.begin(), edges.end(),
         [](auto a, auto b) { return get<2>(a) < get<2>(b); });
    DSU dsu(n);
    long long answer = 0;
    int used = 0;
    for (auto [u, v, w] : edges) {
        if (dsu.unite(u, v)) {
            answer += w;
            ++used;
        }
    }
    return used == n - 1 ? answer : -1; // -1 代表不存在生成树
}
\end{lstlisting}

\subsection{强连通分量、割点和桥：最容易混淆的地方}

\begin{itemize}
  \item 强连通分量只用于有向图；割点和桥通常讨论无向图。看到 \lstinline|low| 不要直接套同一份 Tarjan 代码。
  \item 无向图 Tarjan 处理重边时要跳过“父边编号”，不能写成只跳过 \lstinline|v == fa|。
  \item 桥判断是 \lstinline|low[v] > dfn[u]|；割点的根节点需要单独判断 DFS 子树数是否至少为 2。
  \item 有向图拓扑排序的结果数量小于 $n$，说明有环；这不是“少访问了几个点”，而是题目约束与 DAG 假设冲突。
\end{itemize}

\section{图论回归样例：四组就能抓住大多数板子错误}

\begin{enumerate}
  \item \textbf{单点图：} $n=1,m=0$。检查初始化、答案是否为 0、LCA 是否允许 \lstinline|lca(1,1)|。
  \item \textbf{不连通图：} $n=5$，边为 $(1,2),(2,3),(4,5)$。连通块应为 2，点 1 到点 5 应不可达。
  \item \textbf{重边和自环：} 边为 $(1,2,7),(1,2,3),(2,2,0)$。最短路应选权值 3，拓扑排序不能把自环当成正常边。
  \item \textbf{链和菊花：} 一组是 $1-2-3-\cdots-n$，另一组是 1 连接所有其他点。分别检查递归深度和邻接表遍历是否退化。
\end{enumerate}

\begin{tip}
 图论调试时建议同时打印“点数、边数、可达点数、选中边数、拓扑结果长度”。这些量一旦和理论值不符，比只打印最终答案更容易定位是建图、遍历还是转移错了。
\end{tip}

% ==========================================================
\chapter{动态规划}
% ==========================================================

\section{初始化}

\begin{itemize}
  \item 求最大值 $\Rightarrow$ 初始化 $-\infty$（不是 0，收益可能为负）。
  \item 求最小值 $\Rightarrow$ 初始化 $+\infty$。
  \item 不可达状态（\lstinline|dp[j]==inf|）不要参与转移，否则 \lstinline|inf+cost| 溢出。
  \item "前 $i$ 个的最优" vs "以第 $i$ 个结尾的最优"：定义不同，答案位置不同。
\end{itemize}

\section{遍历顺序}

\begin{itemize}
  \item \textbf{区间 DP}：外层枚举长度 len，内层枚举左端点 L，推出 R。
  \item \textbf{0-1 背包}：容量倒序。\textbf{完全背包}：容量正序。
  \item \textbf{树形 DP}：先递归儿子，再转移父亲。
\end{itemize}

\section{转移遗漏}

\begin{itemize}
  \item LIS 每个元素可以单独成为长度 1，别忘初始化 \lstinline|dp[i]=1|。
  \item 计数 DP 有多个否定条件 $\Rightarrow$ 容斥，注意别重复减。
  \item 滚动数组：每轮开头清空当前层，注意 \lstinline|+=| 和 \lstinline|=| 的区别。
\end{itemize}

\section{边界}

\begin{itemize}
  \item 转移含 \lstinline|dp[i-2]| $\Rightarrow$ 循环从 $i=2$ 或 $i=3$ 开始，手写 base case。
  \item 背包：\lstinline|j - weight[i]| 必须 $\ge 0$。
  \item 答案位置：状态是"以 $i$ 结尾" $\Rightarrow$ 答案是 \lstinline|max(dp[1..n])|。
\end{itemize}

\section{先写状态三句话，再写循环}

DP 最常见的错误不是代码语法，而是状态含义在代码中途偷偷变了。开写前强制写下三句话：

\begin{enumerate}
  \item \textbf{状态：} \lstinline|dp[i][j]| 精确表示什么，已经处理了哪些元素，$i,j$ 是否包含当前位置。
  \item \textbf{转移：} 最后一步是什么，从哪些更小的状态转移过来，是否会重复使用同一个物品或元素。
  \item \textbf{答案：} 是 \lstinline|dp[n][m]|、某一行的最大值，还是所有结尾状态的最大值。
\end{enumerate}

例如“前 $i$ 个数的最长上升子序列”和“以第 $i$ 个数结尾的最长上升子序列”不是同一个状态。前者答案通常看 \lstinline|dp[n]|，后者必须看 \lstinline|max(dp[1..n])|；只要把这两个答案位置混用，就会在答案不以最后一个元素结尾的样例上出错。

\begin{note}
 写完转移后，用 $n=0,1,2$ 手算一遍。若连空前缀、单元素和不可能状态都说不清楚，继续写大模板只会把错误藏得更深。
\end{note}

\section{遍历顺序：是否会偷偷重复使用}

背包的循环方向不是格式问题，而是在决定一个物品能否在同一轮被重复使用：

\begin{lstlisting}
// 0-1 背包：每个物品最多使用一次，容量必须倒序
for (int i = 1; i <= n; ++i) {
    for (int j = capacity; j >= weight[i]; --j) {
        dp[j] = max(dp[j], dp[j - weight[i]] + value[i]);
    }
}

// 完全背包：每个物品可以使用多次，容量必须正序
for (int i = 1; i <= n; ++i) {
    for (int j = weight[i]; j <= capacity; ++j) {
        dp[j] = max(dp[j], dp[j - weight[i]] + value[i]);
    }
}
\end{lstlisting}

区间 DP 通常按区间长度从小到大；树形 DP 通常先算儿子再算父亲；依赖前一层的滚动数组必须先确认当前层是否还会读取旧值。一个非常有效的反例是：只有一个物品、重量刚好能放两次的容量。如果 0-1 背包输出了使用两次的价值，循环方向就错了。

\section{不可达状态：不能把 \lstinline|INF| 当成普通数字}

最小值 DP 用正无穷初始化，最大值 DP 用负无穷初始化，但“无穷”只是标记，不是可以参与加法的真实数：

\begin{lstlisting}
const long long INF = (1LL << 60);
vector<long long> dp(n + 1, INF);
dp[0] = 0;

for (int i = 1; i <= n; ++i) {
    if (dp[i - 1] == INF) continue; // 不可达状态不能继续转移
    dp[i] = min(dp[i], dp[i - 1] + cost[i]);
}
\end{lstlisting}

如果把不可达状态写成 0，求最小值时会错误地让它成为最优；如果直接计算 \lstinline|INF + cost|，可能溢出成负数，之后的 \lstinline|min| 会把这个负数当成答案。最大值 DP 同理，不能让 \lstinline|-INF + value| 参与转移。

\section{滚动数组：清空的是“当前层”}

二维 DP 压成一维后，旧状态和新状态共用同一块内存。最常见的错误是上一轮的值残留，或者把本轮需要的 \lstinline|+=| 写成覆盖：

\begin{lstlisting}
vector<long long> old(m + 1), now(m + 1);
for (int layer = 1; layer <= n; ++layer) {
    fill(now.begin(), now.end(), NEG_INF); // 每轮先清当前层
    for (int j = 0; j <= m; ++j) {
        now[j] = old[j];                 // 不选当前物品
        if (j >= w[layer] && old[j - w[layer]] != NEG_INF) {
            now[j] = max(now[j],
                         old[j - w[layer]] + value[layer]);
        }
    }
    swap(old, now);
}
\end{lstlisting}

计数 DP 的滚动数组还要注意加法顺序：每个来源状态可能贡献多次，不能因为看到一个来源就把目标状态清零。多测时，数组大小、组合数、前缀和和记忆化缓存都要在每组开始重新初始化，不能只改 \lstinline|n|。

\section{答案位置、路径恢复与计数取模}

\begin{itemize}
  \item 求“恰好装满”时，不能把所有容量都当成合法答案；只允许读取 \lstinline|dp[capacity]|。
  \item 求“最多装多少”时，容量不必装满，答案通常是 \lstinline|max(dp[0..capacity])|。
  \item 需要输出方案时，必须在转移发生时记录前驱，或者在最终状态上反向判断；只输出最优值不等于完成题目。
  \item 计数 DP 每次加法和乘法都要及时取模；减法要先加模数，不能依赖 C++ 对负数取模的结果。
  \item 多个方案是否有顺序，要决定循环是“先物品后容量”还是“先容量后物品”；两种顺序可能分别统计组合和排列。
\end{itemize}

\section{DP 最小回归样例}

\begin{enumerate}
  \item \textbf{全负收益：} 数组 $[-5,-2,-7]$，最大值应为 $-2$，可以抓出把最大值 DP 错误初始化为 0。
  \item \textbf{答案不在最后：} 数组 $[3,1,2]$ 的最长上升子序列长度为 2，若只输出最后一个状态，容易错误输出 1。
  \item \textbf{0-1 与完全背包：} 一个物品重量 2、价值 3，容量 4；0-1 背包答案为 3，完全背包答案为 6。
  \item \textbf{不可达容量：} 重量只有 3，查询容量 1 和 2，必须输出题目规定的不可达标记，不能输出 0 或一个溢出的负数。
  \item \textbf{多测残留：} 第一组使用较大的容量并得到答案，第二组只有一个物品且容量为 0，检查答案数组和滚动数组是否真正清空。
  \item \textbf{计数顺序：} 硬币面值 $1,2$、目标和 3，若统计组合答案为 2，若统计排列答案为 3；先确认题目问哪一种。
\end{enumerate}

\begin{tip}
DP 怎么都调不对时：找最小样例，把 DP 表打印出来，手算一张表，对比第一处不同的格子——那里就是 bug。
\end{tip}

% ==========================================================
\chapter{数据结构}
% ==========================================================

\section{并查集}

\begin{itemize}
  \item 合并必须根对根：\lstinline|fa[find(u)] = find(v)|。
  \item 初始化：\lstinline|iota(fa+1, fa+n+1, 1)| 或 \lstinline|for(i=1;i<=n;i++) fa[i]=i|。
  \item 统计连通块：数 \lstinline|find(i)| 的种类，或数 \lstinline|fa[i]==i| 的个数。
\end{itemize}

\section{线段树}

\begin{itemize}
  \item Lazy tag：区间加是累加 \lstinline|tag[ls]+=tag[u]|，区间赋值才是覆盖。
  \item Pushdown：传 tag $\to$ 更新子节点 tree 值 $\to$ 清父 tag。三步缺一不可。
  \item 区间长度：左儿子用 \lstinline|mid-l+1|（节点范围），不是查询边界 L。
  \item 越界返回：最大值返回 $-\infty$，最小值返回 $+\infty$，求和返回 0。
\end{itemize}

\section{二分}

\begin{itemize}
  \item 出现 \lstinline|l = mid| $\Rightarrow$ mid 上取整 \lstinline|mid = (l+r+1)/2|，否则死循环。
  \item 防溢出：\lstinline|mid = l + (r-l)/2|。
  \item 浮点二分：固定循环 100 次，别用 \lstinline|while(r-l>eps)|。
  \item \lstinline|check(mid)| 里的辅助数组每次都要清空。
\end{itemize}

\section{双指针 / 滑窗}

\begin{itemize}
  \item 左指针不能越过右指针：\lstinline|while(l<=r && ...)|。
  \item 最长区间：修复合法后更新答案。最短区间：即将破坏条件前更新。
  \item 分段统计漏最后一段 $\Rightarrow$ 加哨兵 \lstinline|s += '#'| 或 \lstinline|a[n+1] = inf|。
\end{itemize}

\section{树状数组：下标、前缀和与差分}

树状数组只维护“前缀可合并”的信息。最常见的单点加、区间和模板如下，所有下标都假定为 1-based：

\begin{lstlisting}
struct Fenwick {
    int n;
    vector<long long> tree;
    Fenwick(int n) : n(n), tree(n + 1, 0) {}

    void add(int x, long long value) {
        for (; x <= n; x += x & -x) tree[x] += value;
    }

    long long prefix_sum(int x) const {
        long long result = 0;
        for (; x > 0; x -= x & -x) result += tree[x];
        return result;
    }

    long long range_sum(int l, int r) const {
        if (l > r) return 0;
        return prefix_sum(r) - prefix_sum(l - 1);
    }
};
\end{lstlisting}

\begin{warn}
 \lstinline|add(0, value)| 会在 \lstinline|x += x & -x| 后永远停在 0，造成死循环；\lstinline|prefix_sum(n+1)| 可能越界。离散化之后也必须把排名改成从 1 开始。
\end{warn}

\subsection{区间加、单点查：把差分藏进树状数组}

对原数组做差分 $d_i=a_i-a_{i-1}$。给区间 $[l,r]$ 加 $v$，只需 \lstinline|add(l,v)|、\lstinline|add(r+1,-v)|；查询位置 $x$ 时求差分前缀和。测试时必须包含 $r=n$，否则 \lstinline|r+1| 越界的错误永远不会出现。

\section{单调栈与单调队列}

\subsection{下一个更大元素：栈内保存什么}

维护一个从栈底到栈顶单调递减的下标栈。处理 $a_i$ 时，所有小于当前值的元素都找到答案；相等元素是否弹出要由题目要求决定：求“严格更大”通常弹出 \lstinline|<=|，求“更大或相等”通常只弹出 \lstinline|<|。

\begin{lstlisting}
vector<int> next_greater(const vector<int>& a) {
    int n = (int)a.size();
    vector<int> answer(n, -1), st; // st 保存还没有答案的下标
    for (int i = 0; i < n; ++i) {
        while (!st.empty() && a[st.back()] < a[i]) {
            answer[st.back()] = i;
            st.pop_back();
        }
        st.push_back(i);
    }
    return answer;
}
\end{lstlisting}

\subsection{滑动窗口最小值：队首才是答案}

单调队列保存下标，队列中的值从小到大。每次先删掉窗口左侧的下标，再弹出队尾所有不可能成为最小值的元素，最后读队首：

\begin{lstlisting}
vector<long long> window_min(const vector<long long>& a, int k) {
    deque<int> q;
    vector<long long> answer;
    for (int i = 0; i < (int)a.size(); ++i) {
        while (!q.empty() && q.front() <= i - k) q.pop_front();
        while (!q.empty() && a[q.back()] >= a[i]) q.pop_back();
        q.push_back(i);
        if (i >= k - 1) answer.push_back(a[q.front()]);
    }
    return answer;
}
\end{lstlisting}

\begin{note}
 窗口长度 $k=1$、$k=n$、数组全相等、数组单调递增和单调递减，是单调队列最有价值的五组测试。尤其要确认队列存的是下标而不是值，因为只有下标才能判断元素是否过期。
\end{note}

\section{离散化与区间结构}

当数值很大但只关心相对大小时，先复制、排序、去重，再用 \lstinline|lower_bound| 找排名。离散化不是“把值除以一个数”，重复值必须映射到同一个编号；若题目有区间端点和坐标长度，还要判断是否需要加入相邻点之间的空隙。

\begin{lstlisting}
vector<long long> values = all_values;
sort(values.begin(), values.end());
values.erase(unique(values.begin(), values.end()), values.end());

int rank_of(long long x) {
    return int(lower_bound(values.begin(), values.end(), x) - values.begin()) + 1;
}
\end{lstlisting}

\begin{warn}
 只把出现过的端点压成连续编号，适合统计点；如果要计算覆盖长度，$x=1$ 和 $x=2$ 之间的长度不能凭空消失，常需要把端点、端点加 1 或相邻坐标都放入离散化数组。先写清楚“一个编号代表点还是代表小区间”。
\end{warn}

\section{数据结构回归样例}

\begin{enumerate}
  \item \textbf{树状数组：} $n=1$，先加在 1，再查询 $[1,1]$；随后查询空区间 $[2,1]$，答案应为 0。
  \item \textbf{单调结构：} 数组 $[3,3,3]$，分别测试严格大于与大于等于；窗口 $k=1$ 和 $k=n$ 都必须输出完整正确结果。
  \item \textbf{线段树懒标记：} 初始 $[1,1,1,1]$，整段加 5 后查询中间区间，再单点修改并查询整段，专门逼出漏写 \lstinline|push_down| 的问题。
  \item \textbf{离散化：} 输入值 $[10^9,10^9, -10^9, 0]$，确认重复值排名相同，负数不会因为数组下标类型转换而错位。
\end{enumerate}

% ==========================================================
\chapter{博弈论}
% ==========================================================

博弈题先确认三件事：双方是否轮流操作、是否完全信息、是否存在后手无法操作的终止状态。默认下面讨论“双方最优、不能操作者输”的正常玩法；题目若写了反常玩法、和棋或同时操作，不能直接套用。

\section{必胜态与必败态}

把一个局面记为状态 $s$。若当前玩家能走到某个必败态，当前态就是必胜态；若所有后继状态都是必胜态，当前态就是必败态。石子游戏的记忆化模板如下：

\begin{lstlisting}
vector<int> moves;                 // 允许拿走的石子数
vector<int> memo;                  // -1 未计算，0 必胜，1 必败

bool win(int stones) {
    if (stones == 0) return false; // 没有合法操作，当前玩家输
    int& answer = memo[stones];
    if (answer != -1) return answer;
    answer = 0;                    // 先假设所有走法都救不了
    for (int take : moves) {
        if (take <= stones && !win(stones - take)) {
            answer = 1;             // 走到对手必败态
            break;
        }
    }
    return answer;
}
\end{lstlisting}

这里的返回值是“当前局面是否必胜”，不是“谁最终获胜”。最容易错的是把 \lstinline|memo| 按“上一个人是否赢”来理解，或者忘记先给终止状态赋值。若状态不是石子数而是多个变量，先把所有变量编码成一个整数，或用哈希表记录状态。

\section{Bash 游戏：先找模周期}

每次最多拿 $k$ 个石子时，$n$ 个石子的状态满足：

\begin{itemize}
  \item $n \bmod (k+1)=0$ 是必败态；
  \item 其他状态是必胜态，拿走 $n\bmod(k+1)$ 个即可把对手送回必败态。
\end{itemize}

边界一定要包含 $n=0$、$n=k$、$n=k+1$、$n=2(k+1)$。当允许拿走的集合不是连续的 $1$ 到 $k$ 时，Bash 公式通常失效，应先用必胜必败 DP 打表，再观察是否出现周期。

\section{普通 Nim：异或和为零就是必败}

有若干堆石子，每次从一堆拿任意正数。令

\[
X=a_1\mathbin{\mathop{\mathtt{xor}}}a_2\mathbin{\mathop{\mathtt{xor}}}\cdots\mathbin{\mathop{\mathtt{xor}}}a_n.
\]

$X=0$ 时是必败态；$X\ne0$ 时必胜。构造操作不能只输出“存在”，而要找出一堆：

\begin{lstlisting}
long long xo = 0;
for (long long x : a) xo ^= x;

for (int i = 0; i < n; ++i) {
    long long target = a[i] ^ xo;
    if (target < a[i]) {           // 只能减少，不能把这一堆变大
        cout << i << ' ' << a[i] - target << '\\n';
        break;
    }
}
\end{lstlisting}

恢复答案的检查式是 \lstinline|a[i] ^ xo < a[i]|。如果异或之后更大，就不是合法的“拿走石子”。测试 $[0,0]$、$[1]$、$[1,1]$、$[1,2,3]$ 和含有大于 $2^{31}$ 的数，分别检查全零、单堆、必败、异或抵消和整数类型。

\section{Sprague-Grundy 函数与 mex}

对无环有向状态图，定义 $SG(s)$ 为所有后继状态的 Grundy 数集合的 mex，即“不在集合中的最小非负整数”。终止状态的后继集合为空，所以 $SG=0$：

\begin{lstlisting}
vector<vector<int>> graph;
vector<int> sg;

int get_sg(int u) {
    int& result = sg[u];
    if (result != -1) return result;
    vector<int> seen(graph[u].size() + 2, 0);
    for (int v : graph[u]) {
        int value = get_sg(v);
        if (value < (int)seen.size()) seen[value] = 1;
    }
    result = 0;
    while (result < (int)seen.size() && seen[result]) ++result;
    return result;
}
\end{lstlisting}

多个互不影响的子游戏的总 Grundy 数是各自 SG 值的异或。判断总局面时只需检查异或和是否为 0；若要恢复操作，就枚举某一子游戏的合法后继，寻找使总异或变为 0 的那个后继。

\begin{warn}
SG 记忆化搜索要求状态图无环，或者至少要先处理环。对有环图直接递归会无限递归；对“走过的点不能再走”的游戏，状态通常必须包含访问集合，不能只用当前顶点作为记忆化键。
\end{warn}

\section{反常 Nim、阶梯 Nim 与常见变形}

\begin{itemize}
  \item 反常 Nim 的终止规则相反：拿走最后一颗的人输。若所有堆都为 1，奇数堆是必败；否则仍检查异或和是否为 0。
  \item 阶梯 Nim 中，石子只能向更高或更低的相邻阶移动时，先确认题目规则；常见版本的关键量是奇数层堆的异或，而不是所有堆直接异或。
  \item “每次只能拿固定几种数量”通常是 SG 或必胜必败 DP；“可以从任意一堆拿任意数量”才是普通 Nim。
  \item 看到“最后一步得分”“操作次数限制”“有资源消耗”，先检查是否已经不是纯组合游戏，可能要把轮次或资源加入状态。
\end{itemize}

\section{Minimax 极大极小搜索与 Alpha-Beta 剪枝}

棋盘状态较小、双方都能看到完整局面时，可以枚举博弈树。下面模板返回“当前玩家从该状态能得到的最大得分”，其中轮到最大化玩家时取最大值，轮到最小化玩家时取最小值：

\begin{lstlisting}
int minimax(State& state, int depth, int alpha, int beta, bool maximizing) {
    if (depth == 0 || terminal(state)) return evaluate(state);

    if (maximizing) {
        int best = INT_MIN;
        for (Move move : legal_moves(state)) {
            apply(state, move);
            best = max(best, minimax(state, depth - 1,
                                     alpha, beta, false));
            undo(state, move);
            alpha = max(alpha, best);
            if (alpha >= beta) break; // beta 剪枝
        }
        return best;
    } else {
        int best = INT_MAX;
        for (Move move : legal_moves(state)) {
            apply(state, move);
            best = min(best, minimax(state, depth - 1,
                                     alpha, beta, true));
            undo(state, move);
            beta = min(beta, best);
            if (alpha >= beta) break; // alpha 剪枝
        }
        return best;
    }
}
\end{lstlisting}

调用前必须保证 \lstinline|apply| 与 \lstinline|undo| 完全对称；否则深层搜索结束后，父层状态已经被污染。若同一局面会从不同路径到达，可以把局面编码成字符串或整数做哈希去重，但哈希键必须包含轮到谁、剩余深度和所有影响后续决策的信息。

\section{博弈题的现场检查顺序}

\begin{enumerate}
  \item 先写出“不能行动时谁输”，确认是正常玩法还是反常玩法。
  \item 画出 $0,1,2,3$ 个石子的状态，手算前四个必胜必败态。
  \item 状态无环且子游戏独立，优先考虑 SG；棋盘小且需要选择具体走法，考虑 Minimax。
  \item 需要输出操作时，不能只判断胜负，必须验证操作后状态确实满足目标不变量。
  \item 多测时清空 SG、记忆化数组、棋盘哈希和搜索计数器；尤其不要把上一组的周期结论直接带到下一组。
\end{enumerate}

\section{博弈回归样例}

\begin{enumerate}
  \item \textbf{必胜必败 DP：} 允许拿 1 或 3 个，测试 $n=0,1,2,3,4$，答案依次为败、胜、胜、胜、败。
  \item \textbf{Bash：} $k=2$ 时测试 $n=0,1,2,3,4,6$，所有 $n\bmod3=0$ 的状态应为败。
  \item \textbf{普通 Nim：} $[0,0]$、$[1]$、$[1,1]$、$[1,2,3]$，分别检查零堆、单堆、异或为零和多堆抵消。
  \item \textbf{反常 Nim：} 测试全是 1 的偶数堆和奇数堆，再测试含有大于 1 的堆，确认没有把普通 Nim 公式无条件套用。
  \item \textbf{搜索状态：} 棋盘只有一个格子、无合法操作、只有一个合法操作、同一状态两条路径到达，检查终止判断、回滚和状态去重。
\end{enumerate}

\begin{tip}
 博弈题不要一上来背结论。先用小数据把胜负序列写出来：如果序列出现周期，再证明周期；如果是多个独立部分，再验证操作是否真的只改变其中一部分。结论和样例对不上时，优先怀疑游戏规则理解错，而不是先改代码。
\end{tip}

% ==========================================================
\chapter{数学与精度}
% ==========================================================

\section{取模（再强调）}

\begin{itemize}
  \item 减法：\lstinline|(a - b%MOD + MOD) % MOD|。
  \item 除法：乘逆元，不能直接除。
  \item 连乘：步步取模。
\end{itemize}

\section{组合计数}

\begin{itemize}
  \item "至少包含 A 或 B" $\ne$ 包含A的 + 包含B的（重复算了交集）。用容斥或正难则反。
  \item 组合数函数开头：\lstinline|if(m<0 || m>n) return 0|，保证 \lstinline|fact[0]=1|。
  \item 球相同/不同、盒相同/不同、允不允许空——四个条件决定公式。
\end{itemize}

\section{数论}

\begin{itemize}
  \item LCM 防溢出：\lstinline|a / gcd(a,b) * b|（先除后乘）。
  \item 分解质因数：循环结束后 \lstinline|if(n>1)| 说明剩下的是一个大质因数。
  \item 负数整除：C++ 向零取整，数学向下取整，负数时结果不同。
\end{itemize}

\section{浮点}

\begin{itemize}
  \item 整数幂不用 \lstinline|pow|，用快速幂。
  \item 判完全平方：\lstinline|long long r = round(sqrt(n)); if(r*r==n)|。
  \item double 判等：\lstinline|abs(a-b) < 1e-9|。
  \item 分式比大小：交叉相乘 $a \cdot d$ vs $b \cdot c$（注意分母符号和溢出）。
\end{itemize}

% ==========================================================
\chapter{计算几何}
% ==========================================================

\section{精度}

\begin{warn}
\begin{itemize}
  \item 浮点比较必须用 eps：\lstinline|abs(a-b) < 1e-9|，直接 \lstinline|==| 几乎必错。
  \item eps 取多少？一般 $10^{-9}$ 够用，但坐标范围大（$10^9$）时误差会放大，考虑用 $10^{-7}$ 或整数化。
  \item 能用整数就不用浮点：叉积、点积、距离平方都可以保持整数运算。
  \item \lstinline|atan2| 精度有限，角度排序优先用叉积比较。
\end{itemize}
\end{warn}

\section{向量基础}

\begin{itemize}
  \item 叉积 $\vec{AB} \times \vec{AC} = (B-A) \times (C-A)$，正值 C 在 AB 左侧，负值右侧，零共线。
  \item 点积 $\vec{AB} \cdot \vec{AC}$：正为锐角，负为钝角，零为直角。
  \item 向量旋转 $90°$：$(x,y) \to (-y,x)$（逆时针）。
  \item 单位化：除以模长，注意模长为 0 的退化情况。
\end{itemize}

\section{线段与直线}

\begin{itemize}
  \item \textbf{线段相交判定}：跨立实验（两次叉积符号相反）+ 共线时的投影重叠检查。
  \item 共线特判容易漏：两线段共线且有重叠部分也算相交。
  \item 点在线段上：叉积为 0 且点积 $\le 0$（即投影在线段范围内）。
  \item 直线交点：用叉积公式 $t = \frac{(C-A)\times d_2}{d_1 \times d_2}$，分母为 0 说明平行。
\end{itemize}

\section{凸包}

\begin{warn}
\begin{itemize}
  \item Andrew 算法排序后分上下凸壳，注意最后一个点不要重复加入。
  \item 叉积判断用 \lstinline|<=0| 还是 \lstinline|<0|：允许共线点在凸包上用 \lstinline|<|，不允许用 \lstinline|<=|。
  \item 所有点共线时凸包退化成线段，面积为 0，周长要特判。
  \item $n \le 2$ 时直接特判，不要跑凸包算法。
\end{itemize}
\end{warn}

\section{多边形}

\begin{itemize}
  \item 面积（Shoelace 公式）：$S = \frac{1}{2}\left|\sum_{i=0}^{n-1}(x_i y_{i+1} - x_{i+1} y_i)\right|$，下标取模。
  \item 点在多边形内：射线法（奇偶计数）或转角法。射线法注意射线恰好过顶点的边界情况。
  \item 多边形方向：叉积求面积为正 $\Rightarrow$ 逆时针，为负 $\Rightarrow$ 顺时针。题目要求特定方向时先判断再 reverse。
  \item Pick 定理（格点）：$S = I + \frac{B}{2} - 1$，边上格点数 $= \gcd(|\Delta x|, |\Delta y|)$。
\end{itemize}

\section{圆}

\begin{itemize}
  \item 两圆关系：比较圆心距 $d$ 与 $r_1 \pm r_2$，注意内含/外切/相交/内切/外离五种。
  \item 直线与圆交点：圆心到直线距离 vs 半径，距离用叉积除以方向向量模长。
  \item 三点确定圆（外接圆）：解方程组或用垂直平分线交点，注意三点共线无解。
  \item 最小覆盖圆：随机增量法，期望 $O(n)$，注意先 \lstinline|random_shuffle|。
\end{itemize}

\section{旋转卡壳}

\begin{itemize}
  \item 求凸包直径（最远点对）：对踵点单调移动，不要每条边都从头扫。
  \item 叉积比较面积来移动对踵点：\lstinline|cross(hull[i+1]-hull[i], hull[j+1]-hull[j]) > 0| 时 j 前进。
  \item 凸包点数 $\le 2$ 时直接算距离，不要进旋转卡壳。
\end{itemize}

\section{半平面交}

\begin{itemize}
  \item 有向直线：左侧为合法区域，方向一致才能排序。
  \item 按极角排序，平行的保留最左边那条。
  \item 双端队列维护：队首队尾都可能弹出，最后还要用队首检查队尾。
  \item 结果为空 $\Rightarrow$ 无解；结果无界 $\Rightarrow$ 加一个大矩形边界。
\end{itemize}

\section{常见坑总结}

\begin{tip}
\begin{itemize}
  \item 坐标范围 $10^9$ 时叉积会到 $10^{18}$，必须用 \lstinline|long long|。
  \item 三点共线、线段端点重合、多边形退化——这些边界情况是计算几何 WA 的重灾区。
  \item 排序比较函数里用叉积时，共线情况要用距离做二级排序，否则凸包/极角排序会乱。
  \item 输出精度：题目要求保留几位小数就用 \lstinline|fixed << setprecision(k)|，别用默认格式。
\end{itemize}
\end{tip}

% ==========================================================
\chapter{读题错误}
% ==========================================================

\section{数值范围}

\begin{itemize}
  \item 变量可以是负数或 0 吗？$-10^9 \le a_i \le 10^9$ 时别默认正数。
  \item 负数会破坏单调性、双指针的前提、乘法的符号。
\end{itemize}

\section{核心对象}

\begin{itemize}
  \item \textbf{Permutation vs Array}：排列意味着互不相同、值域 $1 \sim N$，可以用位置映射和置换环。
  \item \textbf{简单图}：没写 "no self-loops and multiple edges" 就默认有。
  \item \textbf{有向 vs 无向}：看清 directed/undirected，"关注"是单向，"朋友"是双向。
\end{itemize}

\section{操作规则}

\begin{itemize}
  \item \textbf{同时 vs 依次}：simultaneously $\Rightarrow$ 开新数组存下一状态，统一修改。
  \item \textbf{代价}：是每删一个花 $X$，还是执行一次操作总共花 $X$？
  \item \textbf{Any vs All}：all 要求全部满足，any 只要存在一个。
\end{itemize}

\section{坐标与格式}

\begin{itemize}
  \item 输入是 $(row, col)$ 还是 $(x, y)$？C++ 数组是 \lstinline|grid[row][col]|。
  \item 输入顺序：不一定先 $N$ 后 $M$，看清楚。
  \item 无解输出：回题面 Output 段确认，别凭经验写 $-1$。
\end{itemize}

\end{document}
